\chapter*{K?T LU?N}
\addcontentsline{toc}{chapter}{K?T LU?N}

KẾT LUẬN

Sau quá trình nghiên cứu, thiết kế và triển khai, nhóm em đã hoàn thành Dịch vụ Lưu thông (Circulation Service – N2) cho hệ thống quản lý thư viện theo đúng yêu cầu của học phần. Dịch vụ đáp ứng đầy đủ các chức năng nghiệp vụ cốt lõi bao gồm quản lý mượn sách (tạo phiếu mượn, duyệt/từ chối), quản lý trả sách (xác nhận trả và kiểm tra tình trạng), quản lý gia hạn, tính toán và theo dõi phí phạt, ghi nhận doanh thu và quản lý đánh giá sách. Hai giao diện người dùng được xây dựng riêng biệt cho độc giả và thủ thư, cung cấp trải nghiệm sử dụng trực quan và đầy đủ tính năng. Về mặt kiến trúc, dự án đã vận dụng thành công nhiều nguyên lý và mẫu thiết kế phần mềm quan trọng. Hệ thống áp dụng kiến trúc microservices để phân tách toàn bộ ứng dụng thư viện thành ba dịch vụ độc lập, cho phép ba nhóm phát triển song song trên các công nghệ khác nhau. Trong nội bộ dịch vụ N2, kiến trúc phân tầng kết hợp với mô hình REST API giúp tổ chức mã nguồn rõ ràng, dễ bảo trì và mở rộng. Cơ chế Dependency Injection của ASP.NET Core, cùng với việc áp dụng các nguyên lý SOLID, đảm bảo mã nguồn có tính module hóa cao và khả năng kiểm thử tốt. Cơ chế Published Event Logs thể hiện ứng dụng của kiến trúc hướng sự kiện, cung cấp nền tảng cho việc đồng bộ trạng thái và mở rộng hệ thống trong tương lai. Các thử nghiệm thay đổi tham số đã chứng minh tính linh hoạt của hệ thống. Các chính sách vận hành như giới hạn mượn, mức phí và thời hạn mượn có thể được điều chỉnh dễ dàng thông qua API cấu hình hoặc thay đổi hằng số, đáp ứng nhu cầu thay đổi của thư viện trong thực tế mà không yêu cầu tái cấu trúc mã nguồn. Tuy nhiên, dự án vẫn còn một số hạn chế cần được cải thiện trong tương lai. Cơ chế Published Event Logs hiện tại hoạt động theo mô hình pull (N3 phải chủ động truy vấn để lấy sự kiện) thay vì push qua message broker như RabbitMQ hay Kafka, dẫn đến độ trễ trong việc đồng bộ trạng thái. Việc xác thực token yêu cầu gọi đồng bộ đến N3 cho mỗi request, tạo ra điểm nghẽn tiềm năng khi tải cao – có thể cải thiện bằng cơ chế cache token validation. Bên cạnh đó, hệ thống chưa có cơ chế phát hiện và xử lý lỗi liên dịch vụ một cách toàn diện như Circuit Breaker hay Retry Pattern. Nhóm em đề xuất một số hướng phát triển tiếp theo cho dự án. Thứ nhất, tích hợp message broker (RabbitMQ hoặc Azure Service Bus) để chuyển đổi cơ chế Event Logs từ

pull sang push, cho phép các dịch vụ nhận sự kiện theo thời gian thực. Thứ hai, bổ sung cơ chế health check endpoint và dashboard giám sát (Prometheus + Grafana) để theo dõi trạng thái toàn hệ thống. Thứ ba, triển khai API Gateway tập trung với các tính năng nâng cao như rate limiting, request aggregation và circuit breaker. Cuối cùng, bổ sung unit test và integration test để đảm bảo chất lượng mã nguồn khi hệ thống phát triển. Qua bài tập lớn này, nhóm em đã có cơ hội áp dụng những kiến thức lý thuyết về kiến trúc phần mềm vào một dự án thực tế, từ khâu phân tích yêu cầu, thiết kế kiến trúc cho đến triển khai và kiểm thử. Đây là trải nghiệm quý báu giúp nhóm hiểu sâu sắc hơn về tầm quan trọng của thiết kế kiến trúc trong phát triển phần mềm và sẵn sàng cho các dự án lớn hơn trong tương lai.
